% Generated draft: requires mathematical and source review.
\documentclass[11pt]{article}
\usepackage[margin=25mm]{geometry}
\usepackage{fontspec}
\setmainfont{Latin Modern Roman}
\usepackage{amsmath,amssymb,mathtools,unicode-math}
\setmathfont{Latin Modern Math}
\usepackage{enumitem,needspace,xurl,hyperref}
\setlength{\parindent}{0pt}
\setlength{\parskip}{5pt}
\setlength{\emergencystretch}{3em}
\newcommand{\N}{\mathbb{N}}
\newcommand{\Z}{\mathbb{Z}}
\newcommand{\Q}{\mathbb{Q}}
\newcommand{\R}{\mathbb{R}}
\newcommand{\C}{\mathbb{C}}
\newcommand{\F}{\mathbb{F}}
\newcommand{\A}{\mathbb{A}}
\newcommand{\PP}{\mathbb{P}}
\newcommand{\E}{\mathbb{E}}
\newcommand{\eps}{\varepsilon}
\newcommand{\dd}{\,\mathrm{d}}
\begin{document}
\section*{UnsolvedMath ID: 20003294}
\textbf{Source problem code:} catalogue entry
\subsection{Definitions and mathematical statement}
F. Ledrappier's Question 22 asks to describe invariant ergodic measures for the fast stable foliation $W^s_+$, and says the same may be asked for Anosov diffeomorphisms. Its preceding dynamical setup and the meaning of foliation-invariance and ergodicity are absent locally. Translation invariance, transverse holonomy invariance, invariance under the dynamics, and ergodicity for a joint action need not be identical notions. The aid fixes a precise affine toral model rather than resolving those choices for general foliations.

Let $A\in GL(n,\Z)$ be hyperbolic, inducing an Anosov automorphism of $\mathbb T^n=\R^n/\Z^n$. Let $V_+$ be its invariant fastest stable real subspace, in the exponential-contraction convention, and set
\[
 H=\overline{(V_++\Z^n)/\Z^n}\subseteq\mathbb T^n,
 \qquad Y=\mathbb T^n/H,\qquad\pi:\mathbb T^n\to Y.
\]
The closure is a compact subtorus and $A$ induces an automorphism $\overline A$ of $Y$. The aid classifies \emph{probability} measures invariant under every translation by a vector in $V_+$. Such a measure is $H$-invariant and is uniquely the Haar-fiber lift of $\nu=\pi_*\mu$. Without choosing a section, this lift is specified by
\[
 \int_{\mathbb T^n}\varphi\,d\mu
 =\int_Y\left(\int_H\varphi(x+h)\,dm_H(h)\right)d\nu(y),
 \qquad\pi(x)=y,
\]
where the inner expression is independent of the representative $x$ and $m_H$ is normalized Haar measure.

The reported dynamical refinements are
\[
 A_*\mu=\mu\quad\Longleftrightarrow\quad
 \overline A_*\nu=\nu,
\]
and, for these invariant measures, ergodicity for the joint action generated by $A$ and the $V_+$ translations is equivalent to $\overline A$-ergodicity of $\nu$. Translation-only ergodicity is a different statement: in this lift model it requires the base measure to be a point mass, so it is not silently substituted for joint ergodicity.

The arithmetic criterion is
\[
 H^\perp=V_+^\perp\cap\Z^n
 =\{k\in\Z^n:k\cdot v=0\text{ for every }v\in V_+\}.
\]
Consequently dense fast leaves, equivalently this annihilator being zero, force the unique translation-invariant probability to be Haar on all of $\mathbb T^n$. If $H$ is proper, arbitrary probabilities on $Y$ give distinct translation-invariant lifts, with the further $A$-invariance and joint-ergodicity requirements exactly as above. No analogous classification is asserted for arbitrary nonlinear Anosov fast foliations.

The extraction then enters the separate Section 4.6 on finite-volume locally symmetric spaces $X=\Gamma\backslash G/K$, with $G$ connected semisimple real algebraic, $K$ maximal compact and $\Gamma$ a lattice. It defines flats as totally geodesic zero-curvature submanifolds, describes them using $\Gamma gAK$ for a Cartan subgroup $A$, reports finiteness of compact flats of bounded volume [150], and relates this to totally real fields of fixed degree with bounded regulator. The maximal-flat/split-Cartan and counting-equivalence qualifications are not given; a formula for maximal flats must not be read as describing every lower-dimensional flat or a finiteness theorem for every such embedded submanifold. This is neighboring background, not an additional clause of Question 22, and [150] is not supplied in full.
\subsection{Short English statement}
The aid classifies affine fast-translation-invariant probabilities as Haar measures along subtorus fibers with arbitrary quotient measures. Dynamics invariance and joint ergodicity then reduce to the quotient. Dense fast leaves force Haar uniqueness, while the nonlinear foliation question and the separate flats section remain distinct.
\subsection{Sources}
\begin{itemize}
\item Local UnsolvedMath registry, permanent ID 20003294, source code catalogue entry; \texttt{aim-other-notes.json}, notes[100], source tag question.
\item AIM, \emph{Emerging applications of measure rigidity}, Question 22, F. Ledrappier, followed by Section 4.6: \url{https://aimath.org/WWN/measrigid/measrigid.pdf}. The supplied verification identifies the merged section boundary; this draft does not independently check it.
\item Neighboring [150] concerns compact flats and regulator counting; no full entry is supplied. The affine toral classification has no separate authoritative full citation in the summary.
\item Machine-generated research attempt 1: partial result, open at the research run, rigor labelled ``proved,'' low novelty confidence; registry partially solved. Not peer reviewed; independently verify before citation. Registry updated 14 August 2026.
\end{itemize}
\end{document}
